% !TeX program = xelatex
% 独立任务报告；在项目根目录运行 make intermediate。
\documentclass[UTF8,a4paper,11pt,oneside,fontset=fandol]{ctexart}
\newcommand{\RunningTitle}{第一研究步骤：CEA 与 CCER 标的比较}
\input{intermediate/templates/shared.tex}
\usepackage{amsmath}
\begin{document}
\thispagestyle{plain}
\begin{center}
{\LARGE\bfseries 第一研究步骤任务报告：CEA 与 CCER 候选标的比较\par}
\vspace{0.4em}
{\small 核验日期：2026 年 9 月 30 日}\par
\end{center}

\section{任务、范围与结论}
课程作业要求在全国碳排放配额（CEA）和国家核证自愿减排量（CCER）之间选择碳期货标的，解释同质性、价格波动与市场潜力，并为交易单位奠定依据。工作区流程将此列为第一研究步骤。本报告核查现行权利与账户制度、2025 年同一自然年的官方市场信息，给出候选标的比较和研究性选择；完整逐日数据收集列为后续步骤。详细底稿见\par\noindent\path{research/contract-design/CEA_CCER_UNDERLYING_COMPARISON_2025.md}。

\textbf{研究建议：暂选全国 CEA。}其合规清缴需求明确，2025 年实现成交规模明显高于新登记 CCER；按全国登记体系界定资产也比按多个项目类别界定 CCER 更容易提出统一交割标准。具体配额年度、交割方式、可用余额、期货上市和账户对接均未证实，故这一选择可在后续日度与交割核查后调整。本文不填任何合约数值参数。

\section{制度事实与候选资产范围}
国务院第 775 号令第 30 条规定，1 单位碳排放配额相当于 1 吨二氧化碳当量；第 9、14 条分别规定年度分配和重点排放单位清缴。重点排放单位及符合规定的其他主体可参与全国配额交易，购买的配额可用于清缴。全国登记系统记录是配额归属的最终依据，登记账户记录持有与状态；现货交易还要求实名交易账户与资金账户（R1、R3、R4）。

CCER 是登记项目所产生并经核证登记的减排量，以吨二氧化碳当量计。项目需满足方法学、真实性、唯一性及额外性等条件；登记系统记录权属与注销状态。符合国家规定的主体须开立登记、交易账户。CCER 可依相应规则用于抵销碳配额清缴，用后在其登记系统注销（R2）。2023、2024 年发电行业通知的抵销上限为 5\%，不能作为未来各行业与年度恒定比例（R5）。

\begin{table}[htbp]
\centering\small
\caption{两种资产的期货标的筛选要点}
\begin{tabularx}{\textwidth}{@{}L{0.19\textwidth}XX@{}}
\toprule
维度 & CEA & CCER\\
\midrule
权利来源 & 全国年度配额分配及登记 & 符合方法学的项目减排量核证登记\\
主要需求 & 重点排放单位清缴，可形成配额价格套保需求 & 自愿用途及符合当期条件的有限抵销需求\\
内部差异 & 配额年度、结转资格、账户及可用状态 & 方法学、项目类型、核证批次、用途及状态\\
交割前核查 & 指定年度、登记可用余额、接收主体资格 & 指定批次/项目、可用余额、用途与接收资格\\
\bottomrule
\end{tabularx}
\end{table}

上述交割核查是\textbf{设计建议}，并非已有碳期货规则。2019—2024 年配额结转方案设申请资格、数量上限和冻结安排：旧年度持仓不能不经检查就与新年度视为同质（R6）。登记中的持有、现货可交易和期货可交割是不同状态；两套现货登记和交易系统的运行事实不能证明其已能对接广期所期货交割。

\section{市场数据、计算与解释}
观察窗口固定为 2025 年 1 月 1 日至 12 月 31 日。CEA 数值来自上海环境能源交易所年度行情，包含挂牌协议、大宗协议和单向竞价；243 个市场运行日来自生态环境部年度信息。CCER 数值来自全国自愿减排交易系统年度信息，该系统说明 2025 年首批新登记 CCER 于 3 月 7 日开始交易，全年成交方式为挂牌协议（R7—R9）。故全年总量可作为同一自然年的实现规模，尚不能代表同方式盘口流动性。

\begin{table}[htbp]
\centering\small
\caption{2025 年全国现货市场数据及直接计算}
\begin{tabularx}{\textwidth}{@{}L{0.31\textwidth}XX@{}}
\toprule
指标 & CEA & CCER\\
\midrule
成交量 & 234,597,856 吨 & 约 8,844,100 吨\\
成交额 & 14,629,891,128.06 元 & 约 626,000,000 元\\
市场运行日 & 243 日 & 203 日，自 3 月 7 日起\\
每运行日平均成交量 & 约 965,423 吨 & 约 43,567 吨\\
最大成交类别 & 大宗协议 154,201,694 吨 & 海上风电项目 777.49 万吨\\
\bottomrule
\end{tabularx}
\end{table}

CCER 的年度量和额分别由官方的 884.41 万吨与 6.26 亿元换算，均为四舍五入数。以相同自然年的公布总量计算，
\begin{align*}
\text{成交量比}&=234{,}597{,}856/8{,}844{,}100\approx26.53,\\
\text{成交额比}&=14{,}629{,}891{,}128.06/626{,}000{,}000\approx23.37,\\
\text{运行日平均量}&:\quad 234{,}597{,}856/243\approx965{,}423,\\
&\hspace{5.9em}8{,}844{,}100/203\approx43{,}567\quad\text{吨/日}.
\end{align*}
运行日平均不是有成交日均量，也不衡量订单簿深度。CEA 大宗协议约占全年总量的 65.7\%；若日后采用现货挂牌价格作最终结算价，必须单独检验挂牌价格代表性。CCER 海上风电项目约占总量的 87.9\%，项目集中也需考虑供应与价格风险。这两个比率分别按 $154{,}201{,}694/234{,}597{,}856$ 与 $777.49/884.41$ 计算（R7、R8）。

价格证据仍不足以比较\textbf{收益率波动}：CEA 年报的 97.01 与 50.34 元/吨是综合价格年度高低，CCER 年报的 107.36 与 50.84 元/吨是日成交均价高低，统计量不同。CCER 两种项目的全年成交均价分别为 81.58 和 69.27 元/吨；差额 12.31 元/吨只是类别汇总差，不能解释为同日质量溢价。完整逐日价格及各年度/批次成交尚未清洗，不能据此定保证金、涨跌停板或交易单位（R7、R8）。

\section{判断依据、局限与下一步骤}
暂选 CEA 基于三项判断：第一，清缴产生明确的潜在套保需求；第二，同年成交规模较大；第三，虽然配额年度需要严格筛选，CEA 的差异维度比多项目、多方法学的 CCER 更便于定义一类标准化资产。第三项是设计推论，不是法规认定。拟议标的只含\textbf{全国登记的 CEA}；地方试点配额和 CCER 不作替代交割品。是否限定某一年度须等账户可用量与结转规则核查后决定。

本步骤已取得年度官方数据及法规原文，\textbf{未取得完整同口径逐日表}。下一研究步骤应分别下载或记录两机构的日行情，保留日期、配额年度或项目批次、交易方式、量、额及价格定义，计算有成交日、零成交日、日量分布、价格变化及履约窗口差异；并取得可交割余额、接收账户资格及系统接口的正式证据。若合格年度挂牌交易不连续，或交割账户无法接收拟议资产，本步骤的 CEA 选择及后续交割方案须重审。

\section*{核对的资料}
\small
\noindent R1. 国务院（2024）。《碳排放权交易管理暂行条例》（第 775 号令）。\url{https://www.mee.gov.cn/zcwj/gwywj/202402/t20240205_1065850.shtml}\par
\noindent R2. 生态环境部、国家市场监督管理总局（2023）。《温室气体自愿减排交易管理办法（试行）》（第 31 号令）。\url{https://www.mee.gov.cn/xxgk2018/xxgk/xxgk02/202310/t20231020_1043694.html}\par
\noindent R3. 生态环境部（2021）。《碳排放权登记管理规则（试行）》。\url{https://www.mee.gov.cn/xxgk2018/xxgk/xxgk01/202105/W020210519636656625912.pdf}\par
\noindent R4. 生态环境部（2021）。《碳排放权交易管理规则（试行）》。\url{https://www.mee.gov.cn/xxgk2018/xxgk/xxgk01/202105/W020210519636657102983.pdf}\par
\noindent R5. 生态环境部（2024）。《关于做好 2023、2024 年度发电行业全国碳排放权交易配额分配及清缴相关工作的通知》。\url{https://www.mee.gov.cn/xxgk2018/xxgk/xxgk03/202410/t20241021_1089750.html}\par
\noindent R6. 生态环境部（2024）。《2019—2024 年度碳排放配额结转方案》（配额方案附 5）。\url{https://www.mee.gov.cn/xxgk2018/xxgk/xxgk06/202407/W020240702328244215540.pdf}\par
\noindent R7. 全国温室气体自愿减排交易系统（2025）。《2025 年全国温室气体自愿减排交易市场运行情况一览》。\url{https://www.ccer.com.cn/wcm/ccer/html/2307ptggc1/20260101/090113981.shtml}\par
\noindent R8. 上海环境能源交易所（2025）。《全国碳市场每年综合价格行情及成交信息 20250102—20251231》。\url{https://overview.cneeex.com/c/2025-12-31/497010.shtml}\par
\noindent R9. 生态环境部（2026）。《2025 年全国碳市场平稳有序运行》。\url{https://www.mee.gov.cn/ywgz/ydqhbh/wsqtkz/202601/t20260101_1139528.shtml}\par
\end{document}
